\chapter{Maxwell Theory, Constraints, and Gauge Fixing}
\label{ch:maxwell-theory-constraints-gauge-fixing}

\section{Field Strength and Action}

Let \(A_\mu(x)\) be a real vector-potential representative on four-dimensional
Minkowski spacetime. Gauge-equivalent representatives are related by
\[
  A_\mu(x)\mapsto A_\mu(x)+\partial_\mu\xi(x),
\]
where \(\xi\) is a real scalar function. The local gauge-invariant field
strength is
\[
  F_{\mu\nu}
  =
  \partial_\mu A_\nu-\partial_\nu A_\mu .
\]
The free Maxwell action is
\[
  S[A]
  =
  -\frac14\int\dd^4x\,F_{\mu\nu}F^{\mu\nu}.
\]
Its Euler--Lagrange equation is
\[
  \partial_\mu F^{\mu\nu}=0.
\]

With \(x^0=t\) and spatial indices \(i,j=1,2,3\), the Lagrangian density is
\[
  \mathcal L
  =
  \frac12\sum_{i=1}^3(\partial_0A_i-\partial_iA_0)^2
  -
  \frac14\sum_{i,j=1}^3(\partial_iA_j-\partial_jA_i)^2 .
\]
This expression is the starting point for the Hamiltonian constraint analysis.

\section{Constraints and Gauge Generators}

At a fixed time, regard \(A_\mu(t,\vec x)\) as canonical coordinates. The
momenta conjugate to them are
\[
  \Pi^\mu(\vec x)
  =
  \frac{\partial\mathcal L}
       {\partial(\partial_0A_\mu(t,\vec x))}.
\]
Thus
\[
  \Pi^0(\vec x)=0,
  \qquad
  \Pi^i(\vec x)=\partial_0A_i-\partial_iA_0=F_{0i}.
\]
The equation
\[
  \chi_1(\vec x)=\Pi^0(\vec x)=0
\]
is the primary constraint. Variation with respect to \(A_0\) gives Gauss'
law,
\[
  \chi_2(\vec x)=\partial_i\Pi^i(\vec x)=0,
\]
which is the secondary constraint obtained by preserving \(\chi_1\) under
time evolution.

The equal-time Poisson bracket is
\[
  \{A_\mu(\vec x),\Pi^\nu(\vec y)\}_{\mathrm P}
  =
  \delta_\mu{}^\nu\delta^3(\vec x-\vec y),
\]
with all \(A\)-\(A\) and \(\Pi\)-\(\Pi\) brackets equal to zero. The
constraints obey
\[
  \{\chi_a(\vec x),\chi_b(\vec y)\}_{\mathrm P}=0,
  \qquad
  a,b=1,2.
\]
They are first-class constraints.

For test functions \(\epsilon,\widetilde\epsilon\), define
\[
  G[\epsilon,\widetilde\epsilon]
  =
  \int\dd^3\vec y\,
  \left[
    \epsilon(\vec y)\Pi^0(\vec y)
    +
    \widetilde\epsilon(\vec y)\partial_i\Pi^i(\vec y)
  \right].
\]
The induced infinitesimal change of representatives is
\[
  \delta A_\mu(\vec x)
  =
  \{A_\mu(\vec x),G[\epsilon,\widetilde\epsilon]\}_{\mathrm P},
\]
hence
\[
  \delta A_0(\vec x)=\epsilon(\vec x),
  \qquad
  \delta A_i(\vec x)=-\partial_i\widetilde\epsilon(\vec x).
\]
If a spacetime function \(\zeta\) satisfies
\[
  \zeta(t,\vec x)=-\widetilde\epsilon(\vec x),
  \qquad
  \partial_0\zeta(t,\vec x)=\epsilon(\vec x),
\]
then
\[
  \delta A_\mu(t,\vec x)=\partial_\mu\zeta(t,\vec x).
\]
The first-class constraints therefore generate the infinitesimal gauge
redundancy of the vector-potential representative.

\section{Axial Gauge Phase Space}

A gauge condition chooses one representative from each gauge orbit in a
domain where the choice is regular. Axial gauge takes
\[
  A_3=0.
\]
In this gauge the variables \(A_3,\Pi^3\) are removed. The Lagrangian becomes
\[
\begin{aligned}
  \mathcal L_{\mathrm{ax}}
  &=
  \frac12\sum_{i=1}^2(\partial_0A_i-\partial_iA_0)^2
  +\frac12(\partial_3A_0)^2  \\
  &\quad
  -\frac12(\partial_1A_2-\partial_2A_1)^2
  -\frac12\sum_{i=1}^2(\partial_3A_i)^2 .
\end{aligned}
\]
The equation obtained by varying \(A_0\) is
\[
  \sum_{i=1}^2\partial_i(\partial_iA_0-\partial_0A_i)
  +\partial_3^2A_0=0,
\]
or
\[
  \sum_{i=1}^2\partial_i\Pi^i+\partial_3^2A_0=0.
\]
Together with \(\Pi^0=0\), this is a second-class pair for modes on which
\(\partial_3^2\) is invertible:
\[
  \left\{
    \Pi^0(\vec x),
    \sum_{i=1}^2\partial_i\Pi^i(\vec y)+\partial_3^2A_0(\vec y)
  \right\}_{\mathrm P}
  =
  \partial_3^2\delta^3(\vec x-\vec y).
\]
For nonzero \(k_3\) Fourier modes, solve
\[
  \partial_3^2A_0
  =
  -\sum_{i=1}^2\partial_i\Pi^i
\]
and eliminate \(A_0,\Pi^0\). The remaining canonical variables are
\[
  A_i(\vec x),\Pi^i(\vec x),
  \qquad
  i=1,2.
\]
They are the two canonical degrees of freedom corresponding to the two photon
helicities.

\section{Gauge Slices in Functional Integrals}

The formal expression
\[
  Z=\int DA_\mu\,\ee^{\ii S[A]}
\]
integrates over all representatives of each gauge-equivalence class. A gauge
condition inserts a slice through the representative space.

\begin{center}
\begin{tikzpicture}[>=Stealth,scale=0.92]
  \draw[thick] (0,0) rectangle (6.4,3.4);
  \node[anchor=east] at (6.12,3.12) {representative space};
  \foreach \x in {1.1,2.9,4.7} {
    \draw[blue!70!black,thick]
      plot[smooth] coordinates {
        (\x,0.25) (\x-0.20,1.05) (\x+0.18,2.05) (\x+0.55,3.05)
      };
  }
  \draw[red!75!black,thick]
    plot[smooth] coordinates {
      (0.45,0.80) (1.7,1.02) (3.15,1.17) (4.8,1.19) (5.9,1.08)
    };
  \node[red!75!black] at (5.0,1.55) {gauge slice};
  \foreach \p in {(1.05,0.88),(2.95,1.15),(4.67,1.18)} {
    \fill[green!55!black] \p circle (2.2pt);
  }
  \node[blue!70!black,align=center] at (1.00,2.35) {gauge\\orbit};
\end{tikzpicture}
\end{center}

Let \(\mathcal G(A)\) be a linear gauge condition. In an abelian theory the
Faddeev--Popov determinant for a linear gauge is independent of \(A_\mu\), so
it contributes only to normalization. For axial gauge,
\[
  \mathcal G(A)=A_3,
\]
the gauge-fixed expression is
\[
  Z_{\mathrm{ax}}
  =
  \int DA_\mu\,\delta(A_3)\,\ee^{\ii S[A]}.
\]

After integrating the Maxwell action by parts,
\[
  S[A]
  =
  \frac12\int\dd^4x\,
  A_\mu
  \left(\eta^{\mu\nu}\Box-\partial^\mu\partial^\nu\right)
  A_\nu,
\]
where
\[
  \Box=\partial^\rho\partial_\rho .
\]
The quadratic operator
\[
  D^{\mu\nu}=\eta^{\mu\nu}\Box-\partial^\mu\partial^\nu
\]
has zero directions \(A_\mu=\partial_\mu\xi\). Gauge fixing removes those
directions from the Gaussian integral.

\section{Covariant Gauges}

For perturbative calculations it is useful to preserve manifest Lorentz
covariance. Let
\[
  \mathcal G(A)=\partial_\mu A^\mu .
\]
Instead of imposing \(\mathcal G(A)=0\) sharply, average over
\[
  \partial_\mu A^\mu=\varphi
\]
with a Gaussian weight. In Euclidean signature this insertion is
\[
  \int D\varphi\,
  \exp\left[-\frac{1}{2\xi_{\mathrm g}}
    \int\dd^4x_E\,\varphi^2\right]
  \delta(\partial_\mu A^\mu-\varphi),
  \qquad
  \xi_{\mathrm g}>0.
\]
Performing the \(\varphi\) integral gives the Euclidean gauge-fixing term
\[
  \Delta S_E
  =
  \frac{1}{2\xi_{\mathrm g}}
  \int\dd^4x_E\,(\partial_\mu A^\mu)^2 .
\]
The corresponding Lorentzian gauge-fixed action is
\[
  S_{\xi_{\mathrm g}}[A]
  =
  \int\dd^4x\,
  \left[
    -\frac14F_{\mu\nu}F^{\mu\nu}
    -
    \frac{1}{2\xi_{\mathrm g}}(\partial_\mu A^\mu)^2
  \right].
\]
Equivalently,
\[
  S_{\xi_{\mathrm g}}[A]
  =
  \frac12\int\dd^4x\,
  A_\mu
  \left[
    \eta^{\mu\nu}\Box
    -
    \left(1-\frac1{\xi_{\mathrm g}}\right)
    \partial^\mu\partial^\nu
  \right]
  A_\nu .
\]

In momentum space,
\[
  \widetilde D_{\xi_{\mathrm g}}^{\mu\nu}(k)
  =
  -\eta^{\mu\nu}k^2
  +
  \left(1-\frac1{\xi_{\mathrm g}}\right)k^\mu k^\nu .
\]
Its inverse is
\[
  \bigl(\widetilde D_{\xi_{\mathrm g}}^{-1}(k)\bigr)_{\mu\nu}
  =
  -\frac{\eta_{\mu\nu}}{k^2}
  +
  (1-\xi_{\mathrm g})\frac{k_\mu k_\nu}{(k^2)^2}.
\]
The Lorentzian Feynman boundary value gives
\[
  \langle\Omega|
  \operatorname T\{\widehat A_\mu(x)\widehat A_\nu(y)\}
  |\Omega\rangle_{\xi_{\mathrm g}}
  =
  \int\frac{\dd^4k}{(2\pi)^4}
  \ee^{\ii k\cdot(x-y)}
  \frac{-\ii}{k^2-\ii\epsilon}
  \left[
    \eta_{\mu\nu}
    -
    (1-\xi_{\mathrm g})
    \frac{k_\mu k_\nu}{k^2-\ii\epsilon}
  \right].
\]
The value \(\xi_{\mathrm g}=1\) is Feynman gauge:
\[
  \langle\Omega|
  \operatorname T\{\widehat A_\mu(x)\widehat A_\nu(y)\}
  |\Omega\rangle_{\xi_{\mathrm g}=1}
  =
  \int\frac{\dd^4k}{(2\pi)^4}
  \ee^{\ii k\cdot(x-y)}
  \frac{-\ii\eta_{\mu\nu}}{k^2-\ii\epsilon}.
\]

\section{Field-Strength Correlators}

The field strength is obtained by differentiating the potential. Therefore
the time-ordered two-point function of \(F_{\mu\nu}\) is
\[
\begin{aligned}
  &\langle\Omega|
  \operatorname T\{
    \widehat F_{\mu\nu}(x)\widehat F_{\rho\sigma}(y)
  \}
  |\Omega\rangle  \\
  &=
  \int\frac{\dd^4k}{(2\pi)^4}
  \ee^{\ii k\cdot(x-y)}
  \frac{-\ii}{k^2-\ii\epsilon}
  \bigl(
    k_\mu k_\rho\eta_{\nu\sigma}
    -k_\nu k_\rho\eta_{\mu\sigma}
    -k_\mu k_\sigma\eta_{\nu\rho}
    +k_\nu k_\sigma\eta_{\mu\rho}
  \bigr).
\end{aligned}
\]
All terms proportional to the gauge parameter cancel because antisymmetrizing
\(k_\mu A_\nu-k_\nu A_\mu\) kills every \(k_\mu k_\nu\) contribution.

The Wightman two-point function obtained from the positive-frequency boundary
value has the one-photon spectral form
\[
\begin{aligned}
  &\langle\Omega|
  \widehat F_{\mu\nu}(x)\widehat F_{\rho\sigma}(y)
  |\Omega\rangle  \\
  &=
  \sum_{h=\pm1}
  \int
  \frac{\dd^3\vec k}{(2\pi)^3\,2|\vec k|}
  \langle\Omega|\widehat F_{\mu\nu}(x)|k,h\rangle
  \langle k,h|\widehat F_{\rho\sigma}(y)|\Omega\rangle .
\end{aligned}
\]
With the Chapter 17 polarization representatives,
\[
  \langle k,h|\widehat F_{\mu\nu}(x)|\Omega\rangle
  =
  Z_A^{1/2}\ee^{-\ii k\cdot x}
  \left(
    -\ii k_\mu\epsilon_\nu^{h*}(k)
    +\ii k_\nu\epsilon_\mu^{h*}(k)
  \right).
\]
Comparison with the explicit free two-point function fixes the free Maxwell
normalization
\[
  Z_A=1.
\]

The next chapter couples this gauge field to a charged spinor field. The
interaction will be formulated on representative fields, while local
gauge-invariant observables and external photon states are read through the
field strength and helicity data constructed here.
